% from ega2-1-fr.tex
\oldpage[II]{179}
et par suite on a un $P$-morphisme canonique
$q:\mathbf{V}(\mathcal{L})\to C_P$, le diagramme commutatif
\begin{equation}
\label{II.8.8.2.6-fr}
  \xymatrix{
    \mathbf{V}(\mathcal{L})\ar[r]^{\pi}\ar[d]_{q} & X\ar[d]^{r}\\
    C_P\ar[r] & P
  }
\tag{8.8.2.6}
\end{equation}
identifiant $\mathbf{V}(\mathcal{L})$ au produit $C_P\times_P X$
(\hyperref[II.1.5.2-fr]{1.5.2})~; comme $r$ est une immersion ouverte, il en
est donc de même de $q$ (\hyperref[I.4.3.2-fr]{I, 4.3.2}). En outre, la
restriction de $q$ à $\mathbf{V}(\mathcal{L})-j(X)$ applique ce préschéma dans
$E_P$ par (\hyperref[II.8.5.2-fr]{8.5.2}), et le diagramme
\begin{equation}
\label{II.8.8.2.7-fr}
  \xymatrix{
    X\ar[r]^{j}\ar[d]_{r} & \mathbf{V}(\mathcal{L})\ar[d]^{q}\\
    P\ar[r]_{\varepsilon_P} & C_P
  }
\tag{8.8.2.7}
\end{equation}
est commutatif (cas particulier de
(\hyperref[II.8.5.1.3-fr]{8.5.1.3})). Les assertions de
(\hyperref[II.8.8.2-fr]{8.8.2}) résultent aussitôt de ces faits, en prenant
pour $g$ le morphisme composé $h\circ q$.
\end{proof}

\begin{remark}[8.8.3]
\label{II.8.8.3-fr}
Supposons en outre que $Y$ soit un préschéma \emph{noethérien} et $f$ un
morphisme \emph{propre}. Comme $r$ est alors \emph{propre}
(\hyperref[II.5.4.4-fr]{5.4.4}), donc fermé, et que c'est une immersion ouverte
dominante, $r$ est nécessairement un \emph{isomorphisme}
$X\xrightarrow{\sim}P$. On verra en outre au chapitre III
(\hyperref[III.2.3.5.1-fr]{III, 2.3.5.1}), que $\mathcal{S}$ est alors
nécessairement une $\mathcal{O}_Y$-Algèbre \emph{de type fini}. Alors il en
résulte que $\mathcal{S}^{\natural}$ est une
$\mathcal{S}_0^{\natural}$-Algèbre \emph{de type fini}
(\hyperref[II.8.2.10-fr]{8.2.10, (i)} et
\hyperref[II.8.7.2.7-fr]{8.7.2.7})~; comme $C_P$ est $C$-isomorphe à
$\operatorname{Proj}(\mathcal{S}^{\natural})$
(\hyperref[II.8.7.3-fr]{8.7.3}), on voit que le morphisme $h:C_P\to C$ est
\emph{projectif}~; comme le morphisme $r$ est un isomorphisme, il en est de
même de $q:\mathbf{V}(\mathcal{L})\to C_P$, et on en conclut que le morphisme
$g:\mathbf{V}(\mathcal{L})\to C$ est \emph{projectif}. En outre, comme la
restriction de $h$ à $E_P$ est un isomorphisme sur $E$ et que $q$ est un
isomorphisme, la restriction (\hyperref[II.8.8.2.3-fr]{8.8.2.3}) de $g$ est un
\emph{isomorphisme}
$\mathbf{V}(\mathcal{L})-j(X)\xrightarrow{\sim}E$.

Si on suppose de plus que $\mathcal{L}$ est \emph{très ample} pour $f$, on
verra aussi au chapitre III (\hyperref[III.2.3.5.1-fr]{III, 2.3.5.1}), qu'il
existe un entier $n_0>0$ tel que
$\mathcal{S}_n=\mathcal{S}_1^n$ pour $n\geqslant n_0$. On en conclura par
(\hyperref[II.8.7.7-fr]{8.7.7}) que $\mathbf{V}(\mathcal{L})$ s'identifie au
préschéma $Z$ obtenu en \emph{faisant éclater dans le cône affine $C$ le
sous-préschéma des sommets} (identifié à $Y$), et que la \emph{section nulle}
de $\mathbf{V}(\mathcal{L})$ (identifiée à $Y$) est \emph{l'image réciproque}
du sous-préschéma des sommets $Y$ de $C$. Une partie des résultats précédents
peut en fait s'établir sans hypothèse noethérienne~:
\end{remark}

\begin{corollary}[8.8.4]
\label{II.8.8.4-fr}
Soient $Y$ un préschéma (resp. un schéma quasi-compact), $f:X\to Y$ un
morphisme propre, $\mathcal{L}$ un $\mathcal{O}_X$-Module inversible ample pour
$f$. Alors le morphisme (\hyperref[II.8.8.2.1-fr]{8.8.2.1}) est propre (resp.
projectif) et sa restriction (\hyperref[II.8.8.2.3-fr]{8.8.2.3}) est un
isomorphisme.
\end{corollary}

\begin{proof}
Pour prouver que $g$ est propre, on peut se ramener au cas où $Y$ est affine,
et il suffit donc de considérer le cas où $Y$ est un schéma quasi-compact. Le
même raisonnement que dans (\hyperref[II.8.8.3-fr]{8.8.3}) montre d'abord que
$r$ est un \emph{isomorphisme} $X\xrightarrow{\sim}P$~; par suite $q$ est aussi
un isomorphisme, et comme la restriction de $h$ à $E_P$ est un isomorphisme
$E_P\xrightarrow{\sim}E$, on voit déjà que
(\hyperref[II.8.8.2.3-fr]{8.8.2.3}) est un isomorphisme. Reste à prouver que
$g$ est \emph{projectif}.

Comme $f$ est de type fini par hypothèse, on peut appliquer
(\hyperref[II.3.8.5-fr]{3.8.5}) à l'homomorphisme

